\documentclass[11pt]{article}
\usepackage[a4paper,margin=27mm]{geometry}
\usepackage{amsmath,amssymb,amsthm,mathtools}
\usepackage[T1]{fontenc}
\usepackage{lmodern}
\usepackage[colorlinks=true,allcolors=blue]{hyperref}
\hypersetup{pdfauthor={Ji Ho Bae}}
\usepackage{microtype}
\newtheorem{theorem}{Theorem}
\newtheorem{lemma}[theorem]{Lemma}
\newtheorem{corollary}[theorem]{Corollary}
\newcommand{\Reach}{\operatorname{Reach}}
\newcommand{\Id}{\mathcal J}
\title{The minimum number of deferred-acceptance states in a doubly prime instance}
\author{Ji Ho Bae\\\small LawAll AI}
\date{10 October 2026}

\begin{document}
\maketitle
\begin{abstract}
For every $n\ge2$, the minimum number of reachable serial
deferred-acceptance states among doubly prime strict complete
stable-marriage instances of order $n$ is $3\cdot2^{n-1}-1$.
The lower bound follows by selecting one rejection witness for each
nonfirst proposal. The resulting connected directed acyclic graph has
two incoming neighbors at each nonsource, and every one of its order
ideals is a reachable state. A counting lemma for such graphs supplies
the bound. A rejection-chain family attains it with $2n-1$ proposals.
This proves the conjecture in Ishida's Section~7(A).
\end{abstract}

\section{The question and its conventions}

An instance has $n$ proposers $P$, $n$ receivers $R$, and strict complete
preference lists. A serial deferred-acceptance state is $(h,q)$, where
$q(p)$ counts proposals made by $p$ and $h(r)$ is the proposer held by
$r$, or $\bot$. Initially all counts are zero and all receivers are empty.
A legal step selects an unheld proposer with an untried receiver and
makes its next proposal. The receiver keeps its preferred of the new
proposer and its previous holder. Write $\Reach(I)$ for all states
reached by all finite legal schedules, including the initial and terminal
states.

An \emph{input-independent partition} pairs partitions
$P=P_1\sqcup\cdots\sqcup P_t$ and $R=R_1\sqcup\cdots\sqcup R_t$ with
$|P_j|=|R_j|>0$ for every $j$, such that every member of $P_j$ ranks all
of $R_j$ above every receiver outside $R_j$, and every member of $R_j$
ranks all of $P_j$ above every proposer outside $P_j$. An instance is
\emph{input-prime} if it has no such partition with $t\ge2$.
Number a proposer's events from $1$, and put
\[
 F(q)=\{(p,k):1\le k\le q(p)\},\qquad
 \mathcal F(I)=\{F(q):(h,q)\in\Reach(I)\},\qquad
 E(I)=\bigcup\mathcal F(I).
\]
The change from zero-based event indices in~\cite{ishida} is only a
relabeling. The instance is \emph{behaviorally prime} if no partition
$E(I)=E_1\sqcup E_2$ into nonempty sets satisfies
\[
 \mathcal F(I)=\{X_1\cup X_2:X_j\in\mathcal F(I)|_{E_j},\ j=1,2\},
 \qquad
 \mathcal F(I)|_{E_j}=\{X\cap E_j:X\in\mathcal F(I)\}.
\]
It is \emph{doubly prime} if both properties hold. Let
\[
 m_n=\min\{|\Reach(I)|: I\text{ is doubly prime of order }n\},
 \qquad \min\varnothing=+\infty.
\]

\begin{theorem}\label{thm:main}
For every $n\ge2$, $m_n=3\cdot2^{n-1}-1$. A minimizing instance can
simultaneously have total proposal count $2n-1$.
\end{theorem}

Ishida~\cite[Proposition~21.4 and Section~7(A)]{ishida} proves the lower
bound $2^n+n-1$ and asks for the equality above. The rejection-chain
construction in Section~\ref{sec:sharp} already appears in the
research attempts for catalogue entry C78-352;\footnote{See the
\href{https://github.com/mehmetmars7/Erdosproblems-llm-hunter/blob/00ced55ba1223fb48603e1f15028d753997ee515/attacks/open_problems/economics/gpt_6_astra_ultra/C78-352.tex}{GPT~6 Astra attempt, ``A chain of rejections''}, and the
\href{https://github.com/mehmetmars7/Erdosproblems-llm-hunter/blob/00ced55ba1223fb48603e1f15028d753997ee515/attacks/open_problems/economics/opus_5.5_high/C78-352.tex}{Opus~5.5 attempt, ``The upper bound for every $n$''}.
Both leave the general lower bound unresolved.}
we include its proof to
make attainability explicit. The contribution here is the universal
matching lower bound. In particular, we do not use the experimentally
checked converse relating proposal overlap to behavioral factors
discussed in~\cite[Sections~2 and~4]{ishida}.

\section{Counting ideals by component merges}

For a finite directed acyclic graph $D$, an \emph{ideal} is a vertex set
closed under predecessors. Let $\Id(D)$ be its set of ideals. Connectivity
below refers to the underlying undirected graph.

\begin{lemma}\label{lem:dag}
If $D$ is connected, has $n\ge1$ sources, and every nonsource has at most
two incoming neighbors, then
\[
                         |\Id(D)|\ge3\cdot2^{n-1}-1.
\]
\end{lemma}
\begin{proof}
Choose a topological order placing all sources first. After these $n$
vertices have been inserted, the processed graph has $n$ singleton
components. Every later vertex attaches to one or two existing
components. Exactly $n-1$ vertices merge two distinct components, since
the final graph is connected. Every processed component contains a
source.

For a merging vertex $v$, let $s(v)$ be the number of sources in the two
components being merged. Every ancestor of $v$ lies in those components.
Let $A(v)$ consist of $v$ and all its ancestors. For every subset $T$ of
the $n-s(v)$ sources outside those components, $A(v)\cup T$ is an ideal.
These $2^{n-s(v)}$ ideals are distinct, since their intersection with the
outside sources recovers $T$. Their latest nonsource in the fixed
topological order is $v$. Thus the families for different merging
vertices are disjoint, and none consists solely of sources. All $2^n$
subsets of the sources are themselves ideals.

Record each component merger as a binary node, with the original source
components as leaves. Vertices that do not merge components leave this
record unchanged. The result is a full binary tree $T$ with $n$ leaves;
$s(v)$ is the number of descendant leaves of its internal node $v$.
For a full binary tree $T$ with $\ell$ leaves, define
\[
             C(T)=\sum_{v\text{ internal}}2^{\ell-s(v)}.
\]
For the merger tree $\ell=n$, so $C(T)$ is the total size of the
families above. We claim $C(T)\ge2^{\ell-1}-1$. For a single leaf both
sides are zero. If the root has subtrees $T_a$ and $T_b$ with $a,b\ge1$
leaves, where $a+b=\ell$, then the root contributes $1$, each internal
node $v$ of $T_a$ contributes $2^{\ell-s(v)}=2^b\cdot2^{a-s(v)}$, and
symmetrically for $T_b$. Induction gives
\begin{align*}
 C(T)&=2^bC(T_a)+2^aC(T_b)+1\\
     &\ge2^\ell-2^a-2^b+1
      \ge2^{\ell-1}-1.
\end{align*}
The last inequality follows from $(2^a-2)(2^b-2)\ge0$, equivalently
$2^a+2^b\le2^{\ell-1}+2$. Adding the source-only ideals proves the lemma.
\end{proof}

\section{Rejection witnesses give reachable ideals}

We first record the execution facts needed to apply
Lemma~\ref{lem:dag}. A receiver always holds its favorite among its past
proposers. Thus $q$ determines $h$, and distinct event sets determine
distinct states. Every run has at most $n^2$ steps. Two enabled steps by
different proposers commute and leave each other enabled: another
proposer's action cannot make an unheld proposer held. If they contact
the same receiver, its holder after both steps is the best of the old
holder and the two newcomers.

Consequently the terminal state is unique. For completeness, induct on
$n^2-\sum_p q(p)$: two distinct first steps have a common successor after
commuting; the inductive uniqueness of terminal descendants then
identifies the terminal states of both branches. Every prefix extends to
this terminal state. Denote its counts by $q^*$ and its events by $E$.
Every proposer is held there. Otherwise an unheld proposer must have
exhausted its list, so all $n$ receivers have received proposals and hold
$n$ distinct proposers, contradicting the existence of an unheld one.
In particular, all $n$ first events occur.

Let $H$ be the graph on proposers joining those who contact a common
receiver among the events in $E$. If $H$ were disconnected, events in
different components would address disjoint sets of receivers. Deleting
the other components' steps from a legal prefix preserves legality;
conversely, any projected prefixes can be interleaved because their
receiver updates are disjoint. Every prefix uses only $E$, since it
extends to the unique terminal state. Hence the full event family is
exactly the product of its component projections, each nonempty.
Behavioral primality therefore implies that $H$ is connected.

\begin{lemma}\label{lem:witness}
For a behaviorally prime order-$n$ instance, there is a connected
directed acyclic graph $D$ on $E$ with exactly $n$ sources and exactly two
incoming neighbors at each nonsource, such that every ideal of $D$ is a
distinct reachable state.
\end{lemma}
\begin{proof}
Fix a complete legal execution and write $e(p,k)$ for its $k$th proposal
by $p$. For each $k>1$, let $r$ be the receiver contacted by
$e(p,k-1)$. Immediately before $e(p,k)$, proposer $p$ is unheld and $r$
holds some $w$ it prefers to $p$. Choose $w$'s earlier proposal $e(w,j)$
to $r$ as a witness, and add the edges
\[
       e(p,k-1)\longrightarrow e(p,k),\qquad
       e(w,j)\longrightarrow e(p,k).
\]
Both edges point forward in the fixed execution. Their tails are
distinct because $w\ne p$. The sources are exactly the first events.

Let $X$ be an ideal and execute its events in topological order.
Inductively suppose the preceding events form a legal prefix. The
own-proposer edges ensure that, when $e(p,k)$ is next, $p$ has made
exactly its first $k-1$ proposals. For $k=1$ it is unheld. For $k>1$, the
witness has already proposed to $p$'s previous receiver, whose holder
only improves; that receiver therefore cannot hold $p$. Neither can an
earlier receiver: $p$ was unheld before its following proposal, and
another proposer's action cannot make it held again. Thus $p$ is unheld
and its next proposal is legal. This proves reachability of $X$.
Distinct ideals give distinct proposal counts and hence distinct states.

It remains to prove connectivity. Contract each proposer's own event
chain. Fix a receiver $r$. Every proposer $p$ who contacted $r$ except
its final holder has a subsequent own proposal: otherwise $p$ would be
unheld at termination. Its chosen witness at $r$ is some proposer $w$
strictly preferred to $p$, and the contracted graph contains $p w$.
Iterating such strictly improving witness links among the finite set of
proposers to $r$ reaches its best proposer, the final holder. Therefore
all proposers who contacted $r$ are connected in the contraction. This
connects the endpoints of every edge of $H$. Since $H$ is connected, so
are the contraction and $D$.
\end{proof}

The witness graph need not describe every legal prefix: choosing one
witness selects sufficient prerequisites. Its ideals form a subfamily
of the reachable event family, which is all the lower bound requires.
Lemmas~\ref{lem:dag} and~\ref{lem:witness} prove
\[
       |\Reach(I)|\ge3\cdot2^{n-1}-1
\]
for every behaviorally prime instance, and in particular every doubly
prime instance. The same merger argument gives $|E|\ge2n-1$: besides
the $n$ sources there are $n-1$ distinct merging vertices.

\begin{corollary}\label{cor:joint}
Let $I$ be a behaviorally prime instance of order $n$, and let
$Q^*=\sum_pq^*(p)=|E|$ be its total number of proposals. Then
$Q^*\ge2n-1$ and
\[
       |\Reach(I)|\ge3\cdot2^{n-1}-1+\bigl(Q^*-(2n-1)\bigr).
\]
In particular, $|\Reach(I)|=3\cdot2^{n-1}-1$ implies $Q^*=2n-1$.
\end{corollary}
\begin{proof}
Let $D$ be the graph of Lemma~\ref{lem:witness}; its vertex set is $E$,
and it satisfies the hypotheses of Lemma~\ref{lem:dag}. Fix a
topological order and the merging vertices as in the proof of
Lemma~\ref{lem:dag}. As noted above, the $n$ sources and the $n-1$
merging vertices are distinct, so $Q^*\ge2n-1$; for doubly prime
instances this is also~\cite[Proposition~21.4]{ishida}. Each of the
remaining $Q^*-(2n-1)$ nonsources $u$ merges no components. The set
$A(u)$ consisting of $u$ and its ancestors is an ideal, and its latest
nonsource in the fixed order is $u$, since every ancestor of $u$
precedes $u$. These ideals are therefore distinct. None of them occurs
among the ideals counted in the proof of Lemma~\ref{lem:dag}, each of
which either consists of sources or has a merging vertex as its latest
nonsource. Hence $|\Id(D)|\ge3\cdot2^{n-1}-1+(Q^*-(2n-1))$, and
Lemma~\ref{lem:witness} gives the same bound for $|\Reach(I)|$. If
$|\Reach(I)|=3\cdot2^{n-1}-1$, this bound gives $Q^*\le2n-1$, and
equality follows.
\end{proof}

\section{Attainment and double primality}\label{sec:sharp}

For $n\ge2$, take proposers $p_0,p_1,\ldots,p_{n-1}$ and receivers
$r_1,\ldots,r_n$. Let $p_0$ rank $r_1,\ldots,r_n$ in that order.
For $1\le i<n$, let $p_i$ and $r_i$ rank each other first. Complete
every unspecified list arbitrarily to a strict complete order.

Each $p_i$ with $i>0$ makes one proposal and is held permanently at
$r_i$. Only $p_0$ makes later proposals. At $q(p_0)=k-1$, its $k$th
proposal, for $k\ge2$, requires rejection from $r_{k-1}$ by $p_{k-1}$.
It follows that the reachable pointers are exactly
\[
 \begin{array}{ll}
 q(p_0)=0:&q(p_i)\in\{0,1\}\quad(i>0);\\
 q(p_0)=k,\ 1\le k\le n:&q(p_i)=1\quad(1\le i<k),\qquad
 q(p_i)\in\{0,1\}\quad(k\le i<n).
 \end{array}
\]
For sufficiency, execute the selected first proposals of the $p_i$,
then the first $k$ proposals of $p_0$, stopping immediately after the
last one, even if rejected. Consequently
\[
 |\Reach(I_n)|=2^{n-1}+\sum_{k=1}^{n}2^{n-k}
             =3\cdot2^{n-1}-1.
\]

To check behavioral primality directly, name $p_0$'s events
$a_1,\ldots,a_n$ and $p_i$'s event $b_i$. The displayed pointers say
that the event family is precisely the ideals of
\[
          a_1<a_2<\cdots<a_n,\qquad
          b_i<a_{i+1}\quad(1\le i<n).
\]
A mandatory predecessor $e$ of an event $f$ cannot lie in a different
factor of an event-family product: combine a projection containing $f$
with the empty projection in $e$'s factor to contradict that
requirement. All events occur, and the displayed predecessor edges
connect them all. Thus no nontrivial product exists.

For input primality, suppose the block containing $p_0$ has size $k<n$.
Its receivers must be $r_1,\ldots,r_k$, the first $k$ choices of $p_0$.
Since $r_i$ ranks $p_i$ first, that block must contain each
$p_1,\ldots,p_k$ as well as $p_0$, contradicting its size $k$.
Therefore $I_n$ is doubly prime. This proves attainability and completes
Theorem~\ref{thm:main}. Its terminal counts are $q(p_0)=n$ and
$q(p_i)=1$ for $i>0$, giving $2n-1$ proposals. The same instance thus
attains the proposal lower bound of~\cite[Proposition~21.4]{ishida}
and the state minimum.

Ishida~\cite[Section~7(A)]{ishida} observes that every minimizer found
also attains $Q^*=2n-1$ and asks whether the two minima are always
achieved together. The source does not say whether only this direction
is meant or also its converse, so we answer both. Every doubly prime
instance with $m_n$ states has exactly $2n-1$ proposals, since by
Corollary~\ref{cor:joint} an instance with $Q^*\ge2n$ has at least
$3\cdot2^{n-1}$ states. The converse, that every doubly prime instance
with $2n-1$ proposals has $m_n$ states, fails for $n=3$. Let $p_1$ and
$p_2$ rank $r_1,r_2,r_3$, let $p_3$ rank $r_1,r_3,r_2$, and let every
receiver rank $p_1$ first, $p_2$ second and $p_3$ third. Then $p_1$ is
held at $r_1$ from its first proposal on. Proposers $p_2$ and $p_3$
propose first to $r_1$, then to $r_2$ and $r_3$ respectively. A third
proposal by $p_2$ (by $p_3$) needs an earlier rejection at $r_2$ (at
$r_3$), which only an earlier third proposal by $p_3$ (by $p_2$) can
cause, since $p_1$ never addresses these receivers; hence no third
proposal occurs. Write $a$ for the proposal of $p_1$, $b_1,b_2$ for
those of $p_2$, and $c_1,c_2$ for those of $p_3$. Since $r_1$ holds its
favorite past proposer, $b_2$ can be made exactly when $a$ and $b_1$
have been made, and $c_2$ exactly when $c_1$ and at least one of $a$
and $b_1$ have been made. Hence the reachable event sets are exactly the
subsets of $\{a,b_1,b_2,c_1,c_2\}$ obeying these two conditions; for
sufficiency, make the selected first proposals before the second ones.
There are $2+3+3+6=14$ of them, according as they contain neither of
$a,b_1$, only $a$, only $b_1$, or both. The terminal event set is
$\{a,b_1,b_2,c_1,c_2\}$, so $Q^*=5$. Every receiver ranks $p_1$ first,
so every block $P_j$ of an input-independent partition contains $p_1$;
hence $t=1$ and the instance is input-prime. Since $a$ and $b_1$ are
mandatory predecessors of $b_2$, and $c_1$ of $c_2$, the argument above
shows that an event-family product could only have the factors
$\{a,b_1,b_2\}$ and $\{c_1,c_2\}$. It would then contain $\{c_1,c_2\}$,
the union of the projections of $\varnothing$ and $\{a,c_1,c_2\}$, which
is not reachable. So the instance is doubly prime, with $2n-1=5$
proposals and $14>m_3=11$ states.

\section*{Use of AI tools}
The author used research AI agents that the author built and configured (GPT-6 Astra max,
Claude Opus 5.5 max and GPT-6 Astra Pro) during the proof work, for proof search and
construction, proof checking, translation, Python computation and similar tasks. The final
manuscript was written, and the proofs were checked, by the human author. All responsibility
for this work rests with the author personally.

\begin{thebibliography}{9}
\bibitem{ishida}
Y.~Ishida.
\emph{Prime factorisation of stable-matching instances: uniqueness,
simultaneous products, and an exact census}.
arXiv:2610.05617v1, 2026.
\url{https://arxiv.org/html/2610.05617v1}.
\end{thebibliography}
\end{document}
